
\documentclass{article}
\usepackage{colortbl}
\usepackage{makecell}
\usepackage{multirow}
\usepackage{supertabular}

\begin{document}

\newcounter{utterance}

\centering \large Interaction Transcript for game `cladder', experiment `full\_v1.5\_default', episode 6252 with SLED{-}BSU\_\_Qwen3.5{-}27B{-}sft{-}ep1\_\_prm{-}ep1{-}best\_of\_n.
\vspace{24pt}

{ \footnotesize  \setcounter{utterance}{1}
\setlength{\tabcolsep}{0pt}
\begin{supertabular}{c@{$\;$}|p{.15\linewidth}@{}p{.15\linewidth}p{.15\linewidth}p{.15\linewidth}p{.15\linewidth}p{.15\linewidth}}
    \# & \multicolumn{2}{c}{Player} && \multicolumn{2}{c}{Game Master} \\
    \hline

    \theutterance \stepcounter{utterance}  
    & & & \multicolumn{4}{p{0.6\linewidth}}{
        \cellcolor[rgb]{0.9,0.9,0.9}{
            \makecell[{{p{\linewidth}}}]{
                \texttt{\tiny{[P1$\langle$GM]}}
                \texttt{You are an expert in causal inference. The following question is not a typical commonsense query, but rather a meticulously designed question created by a professor specializing in causal inference, intended to assess the students' mastery of the course content.} \\
\texttt{Imagine a self{-}contained, hypothetical world with only the following conditions, and without any unmentioned factors or causal relationships: Unobserved confounders has a direct effect on maternal smoking status and brown eyes. Cigarette tax has a direct effect on maternal smoking status. Maternal smoking status has a direct effect on brown eyes. Unobserved confounders is unobserved. We know that confounder active and high cigarette tax causes smoking mother. confounder active and smoking mother causes brown eyes. We observed the cigarette tax is high and confounder active. Would the infant has blue eyes if smoking mother instead of nonsmoking mother?} \\
\texttt{Start your answer with 'Yes' or 'No', followed by additional reasoning or evidence to support your explanation.} \\
            }
        }
    }
     \\ \\

    \theutterance \stepcounter{utterance}  
    & \multicolumn{4}{p{0.6\linewidth}}{
        \cellcolor[rgb]{0.561,0.737,0.561}{
            \makecell[{{p{\linewidth}}}]{
                \texttt{\tiny{[P1$\rangle$GM]}}
                \texttt{No.  } \\
\\ 
\texttt{In this hypothetical world, the causal structure is:  } \\
\\ 
\texttt{{-} **Unobserved confounder** → **Maternal smoking**  } \\
\texttt{{-} **Unobserved confounder** → **Brown eyes**  } \\
\texttt{{-} **Cigarette tax** → **Maternal smoking**  } \\
\texttt{{-} **Maternal smoking** → **Brown eyes**  } \\
\\ 
\texttt{We are told:  } \\
\texttt{{-} “Confounder active and high cigarette tax causes smoking mother” → so when the confounder is active and tax is high, the mother smokes.  } \\
\texttt{{-} “Confounder active and smoking mother causes brown eyes” → so when the confounder is active and the mother smokes, the infant has brown eyes.  } \\
\\ 
\texttt{We observe:  } \\
\texttt{{-} Cigarette tax = high  } \\
\texttt{{-} Confounder = active  } \\
\\ 
\texttt{Therefore, the observed facts imply:  } \\
\texttt{{-} Maternal smoking = smoking mother (by the first rule)  } \\
\texttt{{-} Infant’s eye color = brown eyes (by the second rule)  } \\
\\ 
\texttt{The question asks: “Would the infant have blue eyes if smoking mother instead of nonsmoking mother?”  } \\
\texttt{But we are already in the situation where the mother smokes (smoking mother). The counterfactual would be: “If the mother did *not* smoke (nonsmoking mother), would the infant have blue eyes?”  } \\
\\ 
\texttt{To answer that, we need to see what eye color would result if maternal smoking were set to nonsmoking, while keeping the confounder active.  } \\
\\ 
\texttt{From the given rules:  } \\
\texttt{{-} “Confounder active and smoking mother causes brown eyes” does **not** say what happens when confounder is active and mother is *not* smoking.  } \\
\texttt{{-} The only rule linking confounder and eye color is that “Unobserved confounder has a direct effect on brown eyes.” This means the confounder itself can cause brown eyes regardless of smoking.  } \\
\\ 
\texttt{Since the confounder is active, and it directly causes brown eyes, the infant would still have brown eyes even if the mother did not smoke. Therefore, changing maternal smoking from smoking to nonsmoking would **not** change the infant’s eye color from brown to blue.  } \\
\\ 
\texttt{Thus, the answer to “Would the infant have blue eyes if smoking mother instead of nonsmoking mother?” is **No**—the infant would still have brown eyes.} \\
            }
        }
    }
    & & \\ \\

    \theutterance \stepcounter{utterance}  
    & & & \multicolumn{2}{p{0.3\linewidth}}{
        \cellcolor[rgb]{0.9,0.9,0.9}{
            \makecell[{{p{\linewidth}}}]{
                \texttt{\tiny{[GM$|$GM]}}
                \texttt{No} \\
            }
        }
    }
    & & \\ \\

    \theutterance \stepcounter{utterance}  
    & & & \multicolumn{2}{p{0.3\linewidth}}{
        \cellcolor[rgb]{0.9,0.9,0.9}{
            \makecell[{{p{\linewidth}}}]{
                \texttt{\tiny{[GM$|$GM]}}
                \texttt{no} \\
            }
        }
    }
    & & \\ \\

    \theutterance \stepcounter{utterance}  
    & & & \multicolumn{2}{p{0.3\linewidth}}{
        \cellcolor[rgb]{0.9,0.9,0.9}{
            \makecell[{{p{\linewidth}}}]{
                \texttt{\tiny{[GM$|$GM]}}
                \texttt{game\_result = WIN} \\
            }
        }
    }
    & & \\ \\

\end{supertabular}
}

\end{document}
